\begin{tikzpicture}[
  x=1mm,y=1mm,
  font=\figurefont\footnotesize,text=NNInk,
  query/.style={rectangle,draw=NNInput,fill=NNInput!18,
    line width=1.1pt,minimum width=36mm,minimum height=24mm,
    text width=32mm,align=center,inner sep=1.5mm},
  output/.style={rectangle,draw=NNOutput,fill=NNOutput!18,
    line width=1.1pt,minimum width=36mm,minimum height=24mm,
    text width=32mm,align=center,inner sep=1.5mm},
  routes/.style={rectangle,draw=NNHidden,fill=NNHidden!18,
    line width=1.1pt,text width=68mm,minimum height=32mm,
    align=left,inner sep=2mm},
  flow/.style={-{Stealth[length=2mm,width=1.4mm]},draw=NNWire,
    line width=1.2pt,rounded corners=1.2mm},
  panel/.style={draw=NNLine,fill=NNPanel,line width=0.85pt}
]
  \draw[panel] (0,40) rectangle (164,78);
  \draw[panel] (0,0) rectangle (164,38);

  \node[query,minimum width=64mm,text width=60mm,minimum height=9mm]
    (shared) at (82,87) {\bfseries 本次查询需要什么答案？};

  \node[query] (posterior-query) at (21,59)
    {{\bfseries 查询一}\\后验分布／边缘概率\\保留不确定性};
  \node[output] (posterior-output) at (65,59)
    {{\bfseries 输出对象}\\近似分布$q$／伪边缘\\局部高斯近似};
  \node[routes] (posterior-routes) at (126,59)
    {近似族：平均场／结构化VI\\
     共轭计算：VMP\\
     数据与复用：SVI／摊销\\
     局部自由能：Bethe／Kikuchi／TRW\\
     局部近似：Laplace／EP\\
     诊断：支持、边缘与收敛};

  \node[query,dashed] (mpe-query) at (21,19)
    {{\bfseries 查询二}\\完整未观测配置MPE\\寻找最大得分};
  \node[output,dashed] (mpe-output) at (65,19)
    {{\bfseries 输出对象}\\合法完整配置\\得分界／最优性证书};
  \node[routes,dashed] (mpe-routes) at (126,19)
    {可行域：ILP与局部多面体松弛\\
     分解求解：对偶分解\\
     树上界：零温度TRW\\
     特殊精确结构：图割\\
     诊断：紧性、取整与间隙};

  \draw[draw=NNWire,line width=1.2pt,rounded corners=1.2mm]
    (shared.south) -- (82,80) -- (1,80) -- (1,59)
    coordinate (shared-hub);
  \draw[flow] (shared-hub) -- (posterior-query.west);
  \draw[flow,dashed] (shared-hub) |- (mpe-query.west);
  \draw[flow] (posterior-query.east) -- (posterior-output.west);
  \draw[flow] (posterior-output.east) -- (posterior-routes.west);
  \draw[flow,dashed] (mpe-query.east) -- (mpe-output.west);
  \draw[flow,dashed] (mpe-output.east) -- (mpe-routes.west);
\end{tikzpicture}
